% \subsection{comparative bsk / sgp4 propagation}
% * TLEs are tied to real spacecraft, which mean that their change in orbit over time is the result of real perturbations acting on the real spacecraft. HYPSO-1, which is the corresponding satellite for the TLEs used by the Skyfield SGP4 models, is an Earth-observation satellite with a fixed imaging payload. This means its attitude will change over time to 
% * TLEs are estimated from observations, which mean there are intrinsic parameter uncertanties.


% \subsection{deployment velocity acquisition fuel consumption}
% % * The formation appears to be a bounded periodic relative orbit with CPO-like characteristics. The larger along-track amplitude is expected for HCW-type relative orbital motion and does not by itself mean the formation is not a CPO.
% % * I am a little careful not to make too hard claims here since I don't really know different formation types that well: 
% %     Limited to a 2 spacecraft concentric formation due to formation-maintenance implementation constraints. This is because the formation targets a classic orbital element difference setpoint, which cannot be used to achieve a precise constant-separation or circular projected orbit formation discussed in \cite{hill2025hydroswarm}.


% \subsection{Steady-state fuel consumption}
% * The corresponding campagin was not run at the full 60-day duration as planned, but ended prematurely at 24 days due to memory overflow issues in the GNSS-R Formation Flight Simulator. 


% \section{GNSS-R Formation Flight Lifetime Estimate}
% * comment on the fact that the linear fitted m\_acq is negative for the zero differential deployment speed case. 
\hl{Introduction}

\section{Propagation Benchmark Interpretation}
\begin{comment}
* 
\end{comment}

\subsection{Basilisk-Skyfield-SGP4 Divergence Interpretation}

Across the same-spacecraft comparison results, the primary thesis Basilisk configuration shows the smallest divergence from the continuously updating Skyfield-SGP4 baseline. Figure \ref{fig:same_sc_ECI_state_mag_diff} shows that the specialization-project exponential atmosphere configuration reaches approximately three times the position difference of the primary thesis configuration after the 14-day propagation interval, while the re-tuned exponential atmosphere configuration reaches approximately 1.5 times the position difference. Since the main difference between the evaluated Basilisk configurations is the atmospheric density model, the observed ordering can be interpreted primarily through the different drag environments experienced by the spacecraft.

The altitude difference in Figure \ref{fig:same_sc_altitude_difference} supports this interpretation. The primary thesis configuration with the NRLMSISE-00 atmosphere model gradually drops below the Skyfield-SGP4 reference altitude, while both exponential atmosphere configurations drift toward higher relative mean altitude. This indicates that the primary thesis configuration experiences stronger effective drag than the exponential atmosphere configurations over the comparison interval. The re-tuned exponential atmosphere model reduces the altitude divergence relative to the specialization-project configuration, but it does not reproduce the same altitude evolution as the NRLMSISE-00 configuration.

The RTN decomposition in Figure \ref{fig:same_sc_pos_RTN_decomposition} shows that the resulting same-spacecraft position difference is dominated by radial and along-track drift. Because we are effectively considering the relative drift between two deactivated spacecraft affected by different amounts of drag, the relative drift mechanics covered in Section \ref{sec:diff_perturb} apply. Sustained differential drag will accumulate into substantial along-track drift. This is because a sustained altitude difference changes the relative orbital period, which accumulates into along-track separation over time. The exponential atmosphere configurations maintain a higher relative mean altitude than the Skyfield-SGP4 baseline and therefore increasingly lag behind the reference trajectory. This produces the negative along-track difference observed in the RTN decomposition. In contrast, the primary thesis configuration drops to a lower relative mean altitude, corresponding to a shorter relative orbital period, and therefore drifts ahead of the Skyfield-SGP4 baseline in the positive along-track direction.

The radial component in Figure \ref{fig:same_sc_pos_RTN_decomposition} should be interpreted together with the dominant along-track separation. Since the RTN frame is defined at the Skyfield-SGP4 leader position, the radial component represents the direction along its position vector away from the Earth. For large along-track separations, this radial projection is not equivalent to the altitude difference, but rather the circular curvature of the orbit. Thus, a sufficiently large along-track difference correspond to a negative radial deviation for spacecraft in the same approximate near-circular orbit.

The same differential-drag interpretation also applies to the same-formation drift difference shown in Figure \ref{fig:same_form_drift_RTN_decomposition}. All Basilisk configurations predict a larger leader-follower along-track drift than the Skyfield-SGP4 baseline. However, the ordering between the Basilisk configurations is reversed and less separated compared to the same-satellite state comparison. The primary thesis configuration gives a slightly larger same-formation along-track drift difference, while the exponential configurations remain close to it. This indicates that the atmospheric model has a much stronger effect on the absolute Earth-relative propagation of each spacecraft than on the passive leader-follower drift over the same interval. The primary thesis configuration likely diverges more than the others, because its differential drag drift is stronger due to each spacecraft being subject to higher intensity drag. Also, the very small relative divergence between the Basilisk configurations are probably due to the \gls{eci}-fixed spacecraft attitude assumption, where the difference in incident flow in each body-fixed frame are only slightly different due to the initial $1^\circ$ in-plane separation angle. 

\begin{comment}
Discuss around the fact that the thesis propagator configuration has the least divergence across all same-spacecraft difference representation:
    * Across all same-spacecraft comparisons of various Basilisk configurations against the updating-TLE Skyfield-SGP4 comparison baseline, the primary thesis configuration show significantly less divergence than all other Basilisk configurations. Figure \ref{fig:same_sc_ECI_state_mag_diff} show that the divergence distance for the specialization-project untuned exponential atmosphere is approximately three times that of the primary thesis NRLMSISE-00 configuration after a 14-day propagation duration. The tuned exponential atmosphere configuration deviation distance is approximately 1.5 times after the same propagation duration. 
    * Since the only difference between the Basilisk configurations are their respective atmosphere models, it is natural to evaluate the deviations in terms of differences in experienced drag. 
    * Figure \ref{fig:same_sc_altitude_difference} showing the altitude same-satellite differences offers important insight. Here, the primary thesis configuration's mean altitude slightly drops below the reference Skyfield-SGP4 altitude, while the tuned and old Basilisk configurations have increasingly higher relative mean altitudes. This strongly indicates that the primary thesis configuration with its NRLMSISE-00 atmosphere model is subject to more drag than the exponential atmosphere counterparts. 
    * Because we are effectively considering the relative drift between two deactivated (no formation control) satellites affected by different amounts of drag, the relative drift mechanics covered in Section \ref{sec:diff_perturb} apply. Sustained differential drag will accumulate into substantial along-track drift, which is exactly what the RTN component divergence shows in Figure \ref{fig:same_sc_pos_RTN_decomposition}. 
    * This is because The drag-induced relative gain in mean altitude for the tuned and specialization-project Basilisk configurations mean that their relative orbital period is extended, which cause them to increasingly lag behind the Skyfield-SGP4 baseline. 
    * The opposite is true for the primary thesis configuration, where more drag causes its mean relative altitude to decrease, which decreases its relative orbital period. This in turn causes this satellite to lead in front of the Skyfield-SGP4 baseline, thus explaining the positive divergence in the along-track direction. 
    * The radial divergence is a consequence of the along-track separation. Because the radial direction is defined along the Skyfield-SGP4 leader's position vector, any along-track same-satellite position difference in the same orbital plane will manifest as a negative radial deviation. 
    
Same-formation position drift difference:
    * The same differential drag argument can be used to explain the observed divergence in the same-formation position drift differences shown in Figure \ref{fig:same_form_drift_RTN_decomposition}. Here, all configurations give a larger along-track leader-follower position drift than the Skyfield-SGP4 baseline, and the previous divergence order has reversed. The primary thesis configuration likely diverges more than the others, because its differential drag drift is stronger due to each spacecraft being subject to higher intensity drag. Also, the very small relative divergence between the Basilisk configurations are probably due to the \gls{eci}-fixed spacecraft attitude assumption, where the difference in incident flow in each body-fixed frame are only slightly different. 

\end{comment}


\subsection{Updating-TLE SGP4 as a Comparison Baseline}

The interpretation of the comparative propagation results depends on the quality of the continuously updating Skyfield-SGP4 baseline. This baseline is not a true physical ground truth, but an observation-updated reference generated from a sequence of \glspl{tle}. Since the propagated \gls{tle} is replaced regularly throughout the comparison interval, the SGP4 prediction error is not allowed to accumulate over the full 14-day horizon in the same way as it would for a static-\gls{tle} propagation. This makes the updating-\gls{tle} baseline a more useful long-horizon comparison reference than the single-\gls{tle} baseline used in the specialization project~\cite{hystad2025}.

The frequent \gls{tle} updates also explain why no distinct correction discontinuities are visible in the comparison figures. Over short propagation intervals, SGP4 is generally expected to remain sufficiently close to the estimated spacecraft state that the transition between adjacent \glspl{tle} does not dominate the plotted state histories. The resulting baseline is therefore best interpreted as a sequence of short SGP4 predictions tied together by updated orbit estimates, rather than as one uninterrupted long-horizon physical propagation.

However, the \glspl{tle} used by the Skyfield-SGP4 baseline are tied to the real HYPSO-1 spacecraft. The orbital evolution represented by the \gls{tle} sequence therefore includes the effect of real perturbations acting on HYPSO-1, including effects associated with its actual attitude and operational history. Since HYPSO-1 is an Earth-observation satellite with a body-fixed imaging payload, its attitude may vary over time due to science objectives. This means that the effective drag experienced by the real spacecraft may differ from the drag experienced by the simplified Basilisk spacecraft, which is propagated with a fixed attitude in the comparative propagation analysis. As a result, differences between Basilisk and the updating-\gls{tle} Skyfield-SGP4 would still be present with a perfect atmosphere model simply because of different experienced attitude-coupled drag.

The updating-\gls{tle} Skyfield-SGP4 trajectory is therefore a valuable comparison baseline, but it should not be interpreted as exact truth. It provides a practical reference for evaluating whether the primary thesis Basilisk configuration improves over the specialization-project configuration and for quantifying the scale of the remaining propagation divergence. The comparison result should consequently be interpreted as agreement with an observation-updated SGP4 reference, rather than as an absolute validation of the Basilisk propagator.
\begin{comment}
    * The evaluation of the primary thesis Basilisk configuration propagation accuracy depends on the accuracy of the comparative model, the updating-TLE Skyfield-SGP4.
    * Because the TLEs are updated every few hours, its prediction accuracy is again dependent on the TLEs, and its physical interpretation. 
    * One thing to point is that although the TLEs are frequently updated, no visible corrections are visible in the propagator comparison figures. This is because the SGP4 TLE propagation algorithm is widely known to be sufficiently accurate for 1-2 days before noticeably deviation from the real state, depending on the orbit and TLE estimation accuracy. Thus, the SGP4 propagation error does not have enough time to accumulate before the propagated TLE is updated.
    * However, the TLEs are tied to real spacecraft, which mean that their change in orbit over time is the result of real perturbations acting on the spacecraft. HYPSO-1, whose TLEs are used in all the Skyfield-SGP4 propagator, is an Earth-observation satellite with a fixed imaging payload. This means its attitude will change over time to perform science operations, causing its experienced drag to vary with time.
    * This means that, even with a perfect Basilisk atmosphere model, deviations will still appear if their attitudes are  different over time. 
    * Still, because the updating TLEs guarantee that the state error of the comparison baseline is small, it still a valuable comparison baseline, and an improvement from the static-TLE baseline considered in the specialization project in \cite{hystad2025}
\end{comment}



\subsection{Implications for the Lifetime Analysis}
\label{sec:comparison_implication_for_lifetime}

The comparative propagation analysis indicates that the primary thesis Basilisk configuration is the most accurate of the evaluated Basilisk configurations relative to the continuously updating Skyfield-SGP4 baseline. Across the same-spacecraft state comparisons, the NRLMSISE-00 configuration produces substantially smaller divergence than both exponential-atmosphere configurations. This supports the use of the primary thesis configuration as the propagation environment for the GNSS-R formation flight lifetime analysis. However, the primary thesis configuration still accumulates a significant same-satellite difference over the 14-day comparison interval, dominated by radial and along-track drift. 

Perhaps more importantly for the lifetime analysis, is the same-formation position drift difference. Although the major drift difference between the primary thesis configuration and the Skyfield-SGP4 baseline remain unexplained, the increased along-track formation drift will require additional formation-maintenance maneuvers. The resulting lifetime estimations will therefore likely be pessimistic in its predictions. However, the updating-TLE Skyfield-SGP4 baseline also limits how strongly the propagation comparison can be interpreted. Since the baseline is an observation-updated SGP4 reference rather than a precise orbit-determination truth model, the observed Basilisk-Skyfield differences are not direct physical prediction errors. The comparison still provides useful evidence that the thesis configuration improves over the specialization-project configuration, but it does not fully validate the long-horizon Basilisk trajectory.









\section{Formation Fuel Consumption}

\subsection{Acquisition Consumption}
\label{sec:discussion_acq_consumption}

The deployment sensitivity results show summarized in Figure \ref{fig:diff_deploy_curve_fitting}, show that acquisition time is generally positively correlated with differential deployment speed, but that the relationship is not strictly increasing. The most visible exception is the $\Delta v_{\mathrm{dep}}=3.5~\mathrm{m/s}$ case, which acquired the formation after $3.92~\mathrm{h}$, compared with $5.48~\mathrm{h}$ for the $\Delta v_{\mathrm{dep}}=3.0~\mathrm{m/s}$ case and $5.43~\mathrm{h}$ for the $\Delta v_{\mathrm{dep}}=4.0~\mathrm{m/s}$ case. Since $t_{\mathrm{acq}}$ is determined directly from the \gls{oed} tolerance and dwell-time criterion, this non-monotonic behavior indicates that the selected acquisition criterion does not provide a perfectly smooth measure of the acquisition transient across the tested deployment cases.

The \gls{oed} histories shown in Figure \ref{fig:diff_deploy_oed} suggest that the acquisition criterion may be somewhat relaxed for the purpose of distinguishing the end of acquisition from the beginning of maintenance. For example, the $\Delta v_{\mathrm{dep}}=4.0~\mathrm{m/s}$ case performs several additional burns after its acquisition time, which further reduces the semi-major-axis difference toward the center of the tolerance band. Similarly, the eccentricity difference is not fully synchronized across all deployment cases until approximately $7.6~\mathrm{h}$, which is later than the largest recorded acquisition time. The cumulative fuel-consumption time series in Figure \ref{fig:diff_deploy_cum_fuel_consumption} also show several post-acquisition burn events of non-negligible size. These observations suggest that the exact numerical value of $t_{\mathrm{acq}}$ is sensitive to the selected \gls{oed} tolerances and dwell time, although the plots alone do not provide a strict separation between late acquisition burns and early maintenance burns.

The acquisition fuel consumption shows a more uniform dependence on differential deployment speed than the acquisition time. The only from this strictly increasing behavior is in the change from $\Delta v_{\mathrm{dep}}=0.0~\mathrm{m/s}$ to $\Delta v_{\mathrm{dep}}=0.5~\mathrm{m/s}$, where the measured acquisition fuel consumption decreases with differential deployment speed. The slightly lower acquisition fuel consumption for the $0.5~\mathrm{m/s}$ case may indicate that a small non-zero differential deployment speed can reduce the fuel needed to initiate the concentric formation. This makes sense considering this removes the need for an initial burn to separate the leader and follower spacecraft. For the remaining cases, the increasing trend is consistent with the larger velocity correction required to remove the initial deployment-induced relative motion and insert the follower into the target formation.

Because the acquisition time is on the order of a couple of hours for the cases studied in this thesis, the time-horizon is too short for any small differential orbital perturbations to take any noticeable effect, and cannot be considered a major driver for acquisition fuel consumption. Thus, the main driver of acquisition fuel consumption is the initial differential deployment condition and the corrective burns required to reach the desired formation. Modeling this cumulative fuel consumption purely as a function of $\Delta v_{\mathrm{dep}}$ is therefore a reasonable simplification. However, as discussed initially, the largest uncertainty in the extracted acquisition fuel values is likely associated with the acquisition-time definition itself, since changing the \gls{oed} tolerances or dwell-time requirement would shift the point at which cumulative fuel consumption is sampled. 

\begin{comment}
* Discuss the acquisition times. Are all transients gone at this point? Are the OED tolerances correctly tuned?
    - As presented in Figure \ref{fig:diff_deploy_curve_fitting}, the acquisition time is positivly correlated with the differential deployment speed. 
    - however, this relationship is not strictly increasing. The most extreme simulated case for this is the 3.5m/s deployment case which had 3.92hours acquisition time, compared to 5.48 for the 3m/s differential deployment case. This is comparable to the 1.5 and 2 m/s cases with 3.88 hours each. The 3.0m/s case also had the highest acquisition time of all cases, even higher than the highest differential deployment speed case at 5.43 hours. 
    - Since the acquisition time is directly determined by the \gls{oed} tolerances, this effect might be explained by too relaxed parameters. Evidence for this is present in Figure \ref{diff_deploy_oed}, where, for example, the 4m/s case performs several additional burns in order to change its semimajor axis closer to the center of the tolerances after its acquisition time. Additionally, we also see that the eccentricity difference is not fully synchronized for all cases until approximately 7.6 hours, which is greater than the highest acquisition time. 
    - This claim of too loose \gls{oed} tolerances is supported by Figure \ref{fig:diff_deploy_cum_fuel_consumption} that show multiple instances of many relatively large burns after the initial acquisition. However, this is not definitive proof as it is hard to distinguish between formation acquisition and maintenance burns by looking at the cumulative fuel consumption


Discuss the simulated acquisition fuel consumption: 
    * The acquisition fuel consumption trend shown, however, is more uniform. 
    * The only case of non-strict increase in acquisition time is in cases 0.0 (5.6g) to 0.5 m/s (5.39) diff deploy speed. 
    * This result is particularly interesting, because it suggests that it is less fuel intensive to achieve a formation with some non-zero differential deployment speed. This makes sense because it removes the need for the follower to expend fuel to accelerate to away from the leader before it can begin the insertion burns for the concentric formation. 
    * Otherwise, it makes sense that the for the acquisition fuel consumption to be strongly positively correlated with the differential deployment speed, because the follower spacecraft has to expend more fuel to correct for a higher the initial differential deployment speed.  



* What are the drivers for acquisition fuel consumption? (Is it reasonable to model acquisition consumption as a function of only $\Delta v_{\mathrm{dep}}$)
    * Because the acquisition time is on the order of a couple of hours for the cases studied in this thesis, the time-horizon is too short for any small differential orbital perturbations to take any noticeable effect, and cannot be considered a major driver for acquisition fuel consumption. 
    * Therefore, I would argue that modeling the acquisition fuel consumption purely as a function of $\Delta v_{\mathrm{dep}}$ is a fair assumption
    * However, because moving the acquisition time also changes the amount of cumulated consumed fuel before this time, the formation tolerances are probably the greatest source of uncertainty for acquisition fuel consumption. 
\end{comment}


\subsection{Steady-State Maintenance Consumption}
\label{sec:discussion_steady_state_fuel_consumption}

The steady-state fuel consumption result in Figure \ref{fig:steady_state_cum_consumption} show that the steady-state cumulative fuel consumption rate is not perfectly constant, but appears to decrease slightly over the 24-day simulation. This trend is physically reasonable, since a lower remaining propellant mass reduces the spacecraft mass, such that a given orbital correction requires slightly less thrust impulse, assuming that the required formation-maintenance corrections remain approximately the same. This behavior differs from the acquisition phase, where fuel consumption is dominated by a short sequence of corrective burns shortly after deployment. Regardless of how the acquisition time is defined, the deployment sensitivity fuel time series in Figure \ref{fig:diff_deploy_cum_fuel_consumption} clearly show that the burn frequency is strongly reduced after the initial formation insertion period. The remaining fuel consumption is therefore associated with slower formation-maintenance corrections, rather than the removal of the initial deployment-induced relative motion.

This behavior differs from the acquisition phase, where fuel consumption is dominated by a short sequence of corrective burns shortly after deployment. Regardless of how the acquisition time is defined, the deployment sensitivity fuel time series in Figure \ref{fig:diff_deploy_cum_fuel_consumption} clearly show that the burn frequency is strongly reduced after the initial formation insertion period. The remaining fuel consumption is therefore associated with slower formation-maintenance corrections, rather than the removal of the initial deployment-induced relative motion.

The steady-state drift mechanism is not expected to be dominated by large ballistic-coefficient differences between the spacecraft. As described in Section \ref{sec:diff_drag}, the differential atmospheric drag depends on differences in ballistic coefficient, which is affected by the drag coefficient and flow-normal cross-sectional area. In the simulated formation, both spacecraft use the same drag coefficient, and Figure \ref{fig:steady_state_eps_modes_final_day} shows that the leader and follower are synchronized in their operational modes during steady-state operation. Since the operational modes determine the attitude profile, the two spacecraft therefore spend most of the time in similar attitudes. This reduces differences in flow-normal area and gives the leader and follower very similar ballistic coefficients.

Some differential drag is still present because the follower follows an oscillatory relative trajectory around the leader. The radial component in Figure \ref{fig:diff_deploy_rtn_relative_position_decomposition} has an amplitude of approximately $400~\mathrm{m}$, which for small separation distances, correspond to a small periodic altitude difference between the leader and follower. The follower therefore periodically occupies slightly higher and lower atmospheric density regions than the leader. However, this separation is small compared with the orbital altitude, and the formation controller also limits long-term along-track drift by keeping the semimajor axis difference close to zero.

The full \gls{oed} time series in Appendix \ref{fig:diff_dep_all_6_oed} show that the \gls{raan} difference increases over time. This indicates a gradual relative precession of the follower orbit plane with respect to the leader. Since the spacecraft are in a highly inclined low Earth orbit, this behavior is consistent with differential $J_2$ perturbation effects from Earth's non-spherical gravity field. Within the simulated steady-state scenario, this makes differential gravitational perturbations a likely contributor to the long-term formation drift that produces the observed maintenance burns.
\begin{comment}
* Explanation of observed fuel consumption:
    * As shown in figure \ref{fig:steady_state_cum_consumption},the cumulative expended fuel increases at a pretty steady rate over the course of 24 days. 
    * Regardless of how you define the acquisition time, figure \ref{fig:diff_deploy_cum_fuel_consumption} clearly show the pace of corrective burns is heavily reduced in the duration ofter the initial 6 days. 
    * This indicates that other mechanisms are the cause of drifting orbital element difference outside their respective tolerances. 
    * The drift caused by differential drag is greatly reduced in this case because of 2 main reasons. Recall from Section \ref{sec:diff:drag} that the main source of differential atmospheric drag is from the ballistic coefficient, which is dependent on both the drag  coefficient and flow-normal cross-section area. First, the spacecraft is configured with a static drag coefficient, which is the same for both spacecraft. Second, Figure \ref{fig:steady_state_eps_modes_final_day} show that their steady-state operational modes are synchronized, which mean they for the majority of the time occupy the same attitude. This greatly reduced the difference between their flow-normal areas, which in total amounts to a very similar ballistic coefficient for the leader and follower.
    * The follower does however have an oscillating radial component with an amplitude of 400m, as shown in Figure \ref{fig:diff_deploy_rtn_relative_position_decomposition}. With this small separation between leader and follower, this translates well into difference in orbital altitude. Consequentially, the follower dips into slightly thicker and thinner parts of the atmosphere, which periodically subjects it to more and less drag compared to the leader. However, this effect is likely very limited due to the small separation distance considered in this formation flight scenario.
    * Because the formation-maintenance controller actively keeps the semimajor axis difference within small tolerances of zero, the natural relative drift in along-track direction is also limited.
    * The full six orbital elements shown in Appendix \ref{fig:diff_dep_all_6_oed} show that the \gls{raan} difference increases over time. This is precession of the orbital plane around the leader spacecraft is consistent with the effect of $J_2$ differential perturbations, which is likely the biggest source of formation drift considered in this scenario. 
\end{comment}




\section{Lifetime Estimate Reliability}


\subsection{Acquisition Fuel Model Reliability}

The acquisition fuel model used in the lifetime estimate is based on the nine deployment velocity cases shown in Figure \ref{fig:diff_deploy_curve_fitting}. The fitted linear model captures the overall increase in acquisition fuel consumption with differential deployment speed, but it does not reproduce all features of the simulated data. The deviations between the predicted and simulated acquisition fuel consumption are shown in Figure \ref{fig:predicted_vs_simulated_remaining_fuel_comparison} and summarized numerically in Table \ref{tab:gnssr_lifetime_prediction_summary}. Some of non-systematic deviation from the linear model may be explained by the uncertainty related to the determination of the acquisition time, as discussed in Subsection \ref{sec:discussion_acq_consumption}. However, the largest deviation occurs for the zero-deployment-speed case, where the fitted model predicts an acquisition fuel consumption of $-2.56~\mathrm{g}$, compared with the simulated value of $5.60~\mathrm{g}$. This gives a prediction error of $8.16~\mathrm{g}$, which is both a relatively large deviation, and is physically unrealistic since acquisition fuel consumption cannot be negative.

The poor zero-speed prediction is consistent with the non-monotonic behavior between the $\Delta v_{\mathrm{dep}}=0.0~\mathrm{m/s}$ and $\Delta v_{\mathrm{dep}}=0.5~\mathrm{m/s}$ cases discussed in Section \ref{sec:discussion_acq_consumption}. Since the remaining data points follow a more regular increasing trend, the fitted linear model is more representative for the higher deployment-speed cases than for the lowest part of the tested range. This indicates that the acquisition fuel dependence on $\Delta v_{\mathrm{dep}}$ is not fully captured by a single unconstrained linear fit across the complete interval.

A higher-order model could potentially represent the low-speed nonlinearity more accurately, but the available data do not clearly support a single quadratic model over the full deployment-speed range. Such a fit would force a continuously curved relationship, even though the simulated acquisition fuel consumption is approximately linear for most of the tested interval. A piecewise model, with a separate low-speed representation below approximately $1~\mathrm{m/s}$ and a linear model at higher differential deployment speeds, might better represent the observed behavior. However, additional differential deployment speed cases would be required to determine this speculated dynamics more closely.

\begin{comment}
How accurate do I think the model predictions are. Discuss the methodology here (specifically data collection method and curve fitting)
* How well does the remaining fuel model fit the simulation data from the deployment cases? 
    * The linear model that was developed to predict acquisition fuel consumption is only based on the nine samples shown in Figure \ref{fig:diff_deploy_curve_fitting}. 
    * The deviations between the predicted and simulated acquisition fuel consumption for every case in the deployment sensitivity campaign is clearly shown in Figure \ref{fig:predicted_vs_simulated_remaining_fuel_comparison}, and summarized with numerical values in Table \ref{tab:gnssr_lifetime_prediction_summary}. 
    * The deviations vary across the board, with no clear trend. The largest predicted deviation from the simulated acquisition fuel consumption is for the zero deployment speed case, with a difference of 8.16m/s, where the predicted fuel consumption is -2.56m/s. First, it is unsurprising that this differential deployment speed yielded the largest deviation from the predicted model because of the unique case around 0-0.5m/s discussed in Section \ref{sec:discussion_acq_consumption}. Second, a negative predicted acquisition fuel consumption is clearly not realistic.
    * This raises concerns about the predictive accuracy of this linear model. A quadratic fitted model might have been able to capture some of the non-linearities around the 0-0.5m/s differential deployment speed range. However, this would have resulted in a continuously growing acquisition fuel consumption, which does not seem to be the case based on the data points shown in Figure \ref{fig:diff_deploy_curve_fitting}. 
    * Perhaps a more accurate model could be built from a quadratic model in the first 0-1.0 differential deployment speed range, and then transitioning to a linear model for speed differences greater than 1.0.
    * In any case, additional differential deployment speed cases should be run to decipher this behavior further. 
    
\end{comment}



\subsection{Steady-State Fuel Model Reliability}

The steady-state maintenance fuel model is based on the post-acquisition cumulative fuel consumption shown in Figure \ref{fig:steady_state_cum_consumption}. A quadratic model was fitted in order to capture the slight decrease in apparent fuel consumption rate over the completed simulation interval due to mass reduction. Although this model follows the simulated fuel trajectory over the available data interval, its extrapolated behavior in Figure \ref{fig:predicted_lifetime} shows that it is not suitable for predicting fuel depletion over approximately 40 days. The corresponding remaining-fuel model does not reach zero fuel, which makes the quadratic lifetime estimate undefined.

A linear steady-state fuel model was therefore also fitted. The comparison between the predicted and simulated remaining fuel time series in Figure \ref{fig:predicted_vs_simulated_remaining_fuel_comparison} shows that, after accounting for the initial offset caused by acquisition-fuel model deviations, the linear model remains approximately parallel to the simulated post-acquisition fuel time series. This indicates that the linear model captures the average steady-state consumption rate over the short post-acquisition interval covered by the deployment sensitivity campaign cases.

The reliability of the linear steady-state model is limited by the available simulation horizon. The deployment sensitivity simulations only provide a few days of post-acquisition comparison, while the long-horizon steady-state campaign was limited to 24 days. As stated in Section \ref{sec:lifetime_limitations}, the campaign was originally intended for a 90-day simulation horizon, but all longer attempts ended prematurely due to persistent memory overflow issues in the GNSS-R Formation Flight Simulator. Longer completed steady-state simulations would be required to model the fuel rate reduction over time due to decreasing stored propellant mass. 

\begin{comment}
* Steady-state consumption
    * In order to capture the slightly decreasing fuel consumption rate over time, a quadratic steady-state cumulated fuel consumption model was fitted to the post-acquisition consumption data shown in Figure \ref{fig:steady_state_cum_consumption}. Although the seemingly good fit for the simulated interval, the use of this model to predict mission lifetime in Figure \ref{fig:predicted_lifetime}, clearly show that this model fails to make any predictions over approximately 50 days. 
    * Therefore, a linear model was also fitted. The difference between its predicted and simulated steady-state fuel consumption for all differential deployment speed cases is shown in Figure \ref{fig:predicted_vs_simulated_remaining_fuel_comparison}. Ignoring the initial offset due to predicted acquisition fuel consumption deviations, we see that the linear model is approximately parallel to the simulated steady-state consumption. 
    * This is a good sign of the model's predictive capabilities, but the short time-horizon of only 2 days is too short to draw any definitive conclusions. 
    * To better model the effect of reducing mass on fuel consumption, longer horizon simulator runs must be completed. As mentioned in Section \ref{sec:lifetime_limitations}, the campaign was intended to run for 90 days, but all attempts ended prematurely due to a persistent memory overflow issue in the GNSS-R Formation Flight Simulator. Either the data collection must be optimized further, or the simulator must be run on hardware that can support its hunger for RAM.
\end{comment}

\subsection{Interpretation of the Predicted Lifetime}

The linear remaining-fuel model gives a predicted fuel-limited lifetime range of $366.5$-$413.6~\mathrm{days}$ across the evaluated deployment-speed cases, as shown in Figure \ref{fig:predicted_lifetime}. The spread in predicted lifetime is caused by the deployment-speed-dependent acquisition fuel consumption, where larger differential deployment speeds reduce the remaining fuel available for long-term formation maintenance. However, the dominant source of uncertainty in the lifetime estimate is arguably the steady-state maintenance fuel consumption rate. This is because the estimated lifetime is extrapolated over hundreds of days, while the acquisition phase lasts only a few hours. For example, reducing the fitted linear maintenance rate from $1.536~\mathrm{g/day}$ to $1.436~\mathrm{g/day}$ increases the predicted lifetime by approximately $28.4~\mathrm{days}$ when zero acquisition fuel consumption is assumed. This is comparable to the lifetime difference between the $\Delta v_{\mathrm{dep}}=0~\mathrm{m/s}$ and $\Delta v_{\mathrm{dep}}=2.5~\mathrm{m/s}$ cases.

The acquisition fuel model also introduces uncertainty, but its effect on the predicted lifetime is smaller. The largest deviation between fitted and simulated acquisition fuel consumption is $8.16~\mathrm{g}$, which corresponds to a lifetime change of approximately $5.3~\mathrm{days}$ when using the fitted linear steady-state maintenance rate. The lifetime estimate is therefore more sensitive to the steady-state maintenance model than to the acquisition model. This also means that the limited duration of the completed steady-state campaign is a central limitation of the final lifetime prediction.

The linear steady-state model does not capture the observed reduction in fuel consumption rate over the completed long-horizon simulation. Since the spacecraft mass decreases as propellant is consumed, the required propellant mass for equivalent orbital corrections may also decrease over time. This makes the linear steady-state model conservative with respect to this effect. In addition, the comparative propagation analysis showed that the NRLMSISE-00 Basilisk configuration used in the GNSS-R Formation Flight Simulator experienced stronger drag than the updating-TLE Skyfield-SGP4 baseline. Although this effect is reduced in the considered leader-follower formation scenario, as discussed Subsection \ref{sec:discussion_steady_state_fuel_consumption}, the same-formation drift comparison in Section \ref{sec:comparison_implication_for_lifetime} showed larger simulated leader-follower drift in Basilisk than in the Skyfield-SGP4 baseline. This indicates that the steady-state fuel consumption may be overestimated, although the increased formation drift is only partially explained by the available comparison results.

Based on these considerations, the overall predicted lifetime range should be interpreted as a conservative fuel-limited estimate for the selected spacecraft configuration, formation geometry, controller, and disturbance environment. The results indicate that the formation can likely remain fuel-capable for over one year under the simulated assumptions, but the exact depletion time is uncertain. If we assume that the linear steady-state fuel consumption rate were overestimated by a maximum of $10\%$, the predicted lifetime would increase by roughly $45~\mathrm{days}$. Combined with the acquisition fuel model uncertainty due to the sensitivity of $t_{\mathrm{acq}}$ to the selected \gls{oed} tolerances, the estimated lifetime is likely accurate to within $2~\mathrm{months}$ of the true lifetime, heavily biased in the positive direction.

\begin{comment}
Predicted lifetime confidence:
    * The final lifetime estimate using the remaining fuel model gives a lifetime range estimate of 366.5-413.6 days, depending on the differential deployment speed, as seen in the extrapolated linear model presented in Figure \ref{fig:predicted_lifetime}. 
    * Arguably, the largest source of lifetime estimation uncertainty stems from the steady-state fuel consumption because of the long durations involved and comparatively small formation-maintenance consumption rate. 
    * For example, decreasing the estimated steady-state fuel consumption from 1.536 g/day to 1.436 g/day cause an increase in the estimated lifetime of 28.4 days, assuming zero acquisition consumption and a linear model. 
    * To put that into perspective, this corresponds roughly to the difference in estimated lifetime between the 0m/s and the 2.5m/s differential deployment speed cases. 
    * There is also an uncertainty in the lifetime estimate from the acquisition fuel consumption estimate, but with a smaller impact. If the predicted acquisition consumption is wrong by 8.16g, the largest deviance observed, the lifetime estimate would change by 5.3 days, assuming the steady-state fuel consumption. 
    
    * As previously discussed, the linear steady-state cumulative fuel consumption model fails to model the decrease in consumption rate over time from reduced stored propellant mass, which makes the derived model conservative. 
    * In addition to this, the propagator comparison showed that the NRLMSISE-00 propagator configuration used in the GNSS-R Formation Flight Simulator subjected the spacecraft to more drag than the updating-TLE Skyfield-SGP4 comparison baseline. Although this effect is greatly reduced for this scenario as discussed in Section \ref{sec:discussion_steady_state_fuel_consumption}, its effect over this long time horizon is likely not negligible. 
    * This again suggest that the steady-state model over-estimates its consumption.
    * Finally, this is again reinforced by the much larger simulated leader-follower drift compared to the Skyfield-SGP4 baseline, as discussed in Section \ref{sec:comparison_implication_for_lifetime}. However, it is important to note that this increased formation drift can only be partially explained, which means this result has to be "taken with a pinch of salt". 
    * Due to these facts, the overall assessment of the lifetime estimates for the different deployment cases, is that the estimates are likely under-estimated. However, quantifying how pessimistic these predictions are is hard to estimate due to the unknown effect of the decreasing stored fuel mass, and the partially unexplained increased formation drift relative to the Skyfield-SGP4 baseline. However, if we assume the linear steady-state model has an error margin of 10\%, this would give a potential increase in estimated lifetime of 45days. Coupled with the unbiased acquisition fuel consumption uncertainty from the linear model and relaxed acquisition time \gls{oed} tolerances, the "true" lifetime is likely within 2 months of the estimated lifetime. 
\end{comment}




\section{GNSS-R Formation Flight Mission Realism}
\begin{comment}
* The operational constraints implemented to limit the satellites operations to be consistent with a real GNSS-R mission did not prevent any formation corrections. It then raises the question, did these operational constraints actually represent the real day-to-day operations of a GNSS-R formation spacecraft? 
* Does the operational modes sufficiently model the realistic operations of formation flight GNSS-R spacecraft? 
* Discuss around the steady-state mode transitions and battery level. 
* Spacecraft configuration (component placement, mass, size, etc.)
* Does the formation meet the GNSS-R formation flight requirements for cooperative measurements? If no, how would maintenance of the preferred formation have changed the observed fuel consumption
\end{comment}

\subsection{Representativeness of the Spacecraft Model}

\subsection{Formation \& Operational Mode Realism}

\subsection{Realism of the Fuel-Consumption Results}
\begin{comment}
    
\end{comment}



\section{Main Limitations and Future Work}
\begin{comment}
* Only one orbital "layer" of the concentric formation was studied. Spacecraft in larger radius relative orbit about the leader will probably experience larger differential perturbations, but this needs more studying. 
* GNSS-R observation geometry
* Fuel consumption related to orbital maintenance is not included. 
\end{comment}